% SHIELD1 OR-BIDIRECTIONAL-1, 2026-10-15 UTC
\documentclass{article}
\usepackage{amsmath,amssymb,amsthm}
\newtheorem{theorem}{Theorem}
\title{SHIELD1: Two-way Resolution and PC Bounds for the OR Shield}
\date{15 October 2026}
\begin{document}\maketitle
\noindent Status: exact mathematical consequences of the OR3 gadget, independently verified in finite cases; no arbitrary SAT decider lower bound.
Let $A,B$ be UNSAT proper three-CNFs on disjoint variables with $m_A,m_B$ clauses and $q=m_A+m_B$. Write
\[
{\cal O}(A,B)=\bigwedge_{C=(a\vee b\vee c)\in A}[(s\vee a\vee y_C)\wedge(\neg y_C\vee b\vee c)]
\wedge\bigwedge_{D=(d\vee e\vee f)\in B}[(\neg s\vee d\vee z_D)\wedge(\neg z_D\vee e\vee f)].
\]
The predecessor SHIELD1 report proves exact existential OR semantics, properness, clause distinctness, $2q$ clauses, $v_A+v_B+q+1$ variables, and literal projections onto both leaves.

\begin{theorem}[Resolution upper and lower bounds]
Let $I(F)$ denote minimal numbers of genuine Resolution inferences, excluding input axiom lines; $W(F)$ is minimal possible maximum width. Then
\[
\max\{I(A),I(B)\}\le I({\cal O})\le I(A)+I(B)+q+1,
\quad
\max\{W(A),W(B)\}\le W({\cal O})\le\max\{W(A),W(B)\}+1.
\]
If input axioms are also displayed, a constructed proof uses at most $3q+I(A)+I(B)+1$ lines.
\end{theorem}
\begin{proof}
Resolve each local pair on its private variable to derive $s\vee C$ for every $A$ clause and $\neg s\vee D$ for every $B$ clause, costing $q$ inferences. Append the same guard to every line of an $A$ refutation and replay its steps. The result is the unit $s$. Likewise the guarded $B$ refutation yields $\neg s$. Resolve these units. Guarding increases every width by one; original clauses had width three and derived guards width four. For lower bounds project any target refutation through the predecessor's literal/constant substitutions, using closure of general Resolution under substitutions. This does not assert regular-Resolution closure.
\end{proof}

\begin{theorem}[PC degree, all fields]
With conventional PC clause polynomials $p_C=\prod_{\ell\in C}(1-\ell)$ and Boolean axioms $x^2-x$, over every field $\mathbb F$,
\[
\max\{D_{\mathbb F}(A),D_{\mathbb F}(B)\}
\le D_{\mathbb F}({\cal O})
\le\max\{D_{\mathbb F}(A)+1,D_{\mathbb F}(B)+1,5\}.
\]
The forward expanded-monomial proof size obeys $S_{\mathbb F}({\cal O})=O(S_{\mathbb F}(A)+S_{\mathbb F}(B)+q)$. This is not a reverse size bound for ordinary PC.
\end{theorem}
\begin{proof}
For the $A$ split put $u=(1-s)(1-a)$, $v=(1-b)(1-c)$, $P=u(1-y)$ and $Q=yv$. The formal polynomial identity
\[
\boxed{vP+uQ=uv=(1-s)p_C}
\]
derives every guarded source clause in PC degree at most five using a constant number of standard polynomial-calculus inferences. The same identity on the $B$ split, with $u=s(1-d)$, derives $sp_D$. Guard each line $R$ of an $A$ source refutation by $1-s$, and each $B$ line by $s$. Original PC inferences commute with these guards; guarded Boolean axioms follow from target Boolean axioms by one variable multiplication and subtraction. The final source constants $1$ become $1-s$ and $s$, whose sum is $1$. Guarding increases source degrees by one and doubles expanded monomial counts at most. The lower bounds follow from degree-preserving literal/constant projection established in SHIELD1.
\end{proof}

\begin{theorem}[Balanced four-leaf portfolio]
For four UNSAT sources $F_i$, $q=\sum_i m_i$, and $\Phi={\cal O}({\cal O}(F_1,F_2),{\cal O}(F_3,F_4))$,
\[
\max_i I(F_i)\le I(\Phi)\le \sum_i I(F_i)+3q+3,\qquad
\max_i W(F_i)\le W(\Phi)\le\max_i W(F_i)+2,
\]
\[
\boxed{\max_i D_{\mathbb F}(F_i)\le D_{\mathbb F}(\Phi)\le\max_i D_{\mathbb F}(F_i)+4}.
\]
\end{theorem}
\begin{proof}
Apply the one-gadget upper constructions at the two lower nodes (combined clause count $q$) and root (child clause count $2q$). Compose the leaf projections for the lower inequalities.
\end{proof}

\noindent Binary semantic-CP also has an additive forward line bound. For any source line $a\cdot x\ge b$, set $M=\max(0,b-\sum_{a_j<0}a_j)$ and guard it by $Ms$ on the $A$ side, or by $M(1-s)$ on the $B$ side. Each valid binary inference remains valid: when its guard is disabled it is the source inference, and when enabled the conclusion is Boolean-valid. Summing the two local split inequalities derives the guarded initial clause. The final guarded contradictions are $s\ge1$ and $1-s\ge1$, semantically implying $0\ge1$. This does not assert an ordinary syntactic-CP forward simulation.

\noindent Verification: a separate constructive Python audit generated 27-line source refutations of two 8-clause full-cube contradictions and a valid 87-line OR-shield Resolution refutation; it also verified 16 exact integer polynomial identities across signs and branch sides, and 28,329 guarded binary semantic-CP inferences. The new result is a no-amplification theorem in Resolution length/width and PC degree, not a computational time lower bound.

\begin{thebibliography}{2}
\bibitem{fhl} Filmus, Hrube\v{s}, Lauria. Semantic Versus Syntactic Cutting Planes, STACS 2016.
\bibitem{prior} SHIELD1 OR3-PROJECTION-AUDIT1, project repository, 15 October 2026.
\end{thebibliography}
\end{document}
% SHIELD1::OR-BIDIRECTIONAL-1::2026-10-15
